\chapter*{Abstract}
\addcontentsline{toc}{chapter}{Abstract}

EduVerse is presented in this thesis as an integrated educational platform
designed to reduce the fragmentation that commonly arises when class
communication, learning materials, live meetings, whiteboard collaboration,
programming practice, assessment, and academic follow-up are distributed across
separate tools. The platform follows a multi-organization, role-based model in
which Organization Administrators, Teachers, and Students operate within a
governed workspace that can adapt to different institutional needs. Within that
workspace, EduVerse supports classes, materials, chat, assignments, exams,
results, notifications, live sessions, whiteboard collaboration, an AI Agent,
an integrated browser-based IDE, configurable feature availability, extension
support, and public or open-learning access where enabled.

The main implementation is a web platform built with Next.js, React, and
TypeScript. It uses Supabase for authentication and relational data, AWS S3 for
file storage, LiveKit for live-session participation, and OpenRouter for
AI-supported learning services. EduVerse also includes a mobile companion
application implemented with Expo, React Native, and TypeScript. The current
mobile client provides authenticated access, class summaries, notifications,
materials, class chat, organization switching, and basic live-session
participation. However, it remains under active development and is not a full
replacement for the web platform, where broader administrative and advanced
learning workflows are still concentrated.

The thesis documents the problem context, system analysis, design,
implementation, and evaluation of EduVerse using the available implementation
evidence, project documentation, and targeted validation. The evaluation shows
credible support for role-based academic workflows, shared API reuse, modular
feature enablement, and integrated educational services, while also identifying
realistic limitations in end-to-end testing breadth and in the current scope of
the mobile companion. Overall, EduVerse contributes a coherent and extensible
approach to building a governed digital learning environment around one
connected academic platform rather than a collection of disconnected tools.
